% [原文 p.176]
可以把 $\nu^{\lambda_\alpha}$ 的结构群约化到群 $SO_{\lambda_\alpha}$），则有不等式
\begin{equation}
  \sum_{\alpha}b_{i-\lambda_\alpha}(W^k_\alpha)\ge b_i(M^n) \tag{2.13}
\end{equation}
其中 $b_j$ 是第 $j$ 个有理同调群的秩（在 $\Z/p$ 上也成立）。若向量丛 $\nu^{\lambda_\alpha}\to W^k_\alpha$ 不全可定向，则不等式 (2.13) 只对 $\Z/2$ 成立。这些不等式可这样解释：临界子流形 $W^k_\alpha$ 中每个 $s$ 维闭链，对应各群 $H_{s+\lambda_\alpha}(M_c,M_{c-\eps})$ 中一个维数为 $s+\lambda_\alpha$ 的相对闭链（这里 $M_c$ 表示由 $f\le c$ 确定的 $M^n$ 的闭区域）；粗略地说，它由从 $s$-闭链的各点 $x$ 出发、沿 $\R^{\lambda_\alpha}_x$ 中「最速下降」方向发出的射线构成。

一般说来，关于流形上可定义的实值光滑函数的分析或几何信息，若对临界点的性质加上种种限制，就能给出关于流形拓扑的重要结论。

\begin{example*}{1}
若闭流形 $M^n$ 上存在只有 2 个临界点的 Morse 函数，则不难看出 $M^n$ 必同胚（甚至分段线性同胚）于 $S^n$。由 Smale--Wallace 定理（20 世纪 60 年代初），对 $n\ge5$，每个同伦等价于球面 $S^n$ 的流形 $M^n$ 上都存在恰有 2 个临界点的 Morse 函数。这种球面刻画的想法属于 Reeb（20 世纪 50 年代初）；它的第一个深刻应用由 Milnor 在 20 世纪 50 年代末给出，与他发现球面 $S^7$ 上的异种光滑结构有关：具有这些光滑结构的流形不与标准光滑 7-球面微分同胚，尽管它们与它 PL 同胚。下面 §4 中我们将讨论关于 $n\ge5$ 时同伦等价于 $S^n$ 的流形 $M^n$ 的、更为精细的 Milnor--Kervaire 理论。
\end{example*}

\begin{example*}{2}
设 $W^{n+1}$ 是带边流形，其边界为不交并 $\partial W^{n+1}=V^n_1\cup V^n_2$，并设 $W^{n+1}$ 可缩到边界的这两部分上：
\[
  W^{n+1}\sim V^n_1\sim V^n_2,\qquad \pi_j(W^{n+1},V^n_i)=0,\quad j\ge0.
\]
这时我们说流形 $W^{n+1}$ 是 $V^n_1$ 与 $V^n_2$ 之间的 \emph{$h$-配边}，并说后两个流形 \emph{$h$-配边}。Smale 的「$h$-配边定理」断言：若 $n\ge5$ 且 $V^n_1$、$V^n_2$ 单连通，则 $W^{n+1}$ 上存在光滑实值函数 $f$，它在 $V^n_1$ 与 $V^n_2$ 上为常数，没有临界点（且其梯度处处不与 $V^n_1$、$V^n_2$ 相切）。由此可知，在同样假设下 $W^{n+1}$ 微分同胚于柱面 $V^n_1\times I$，且 $V^n_1$ 微分同胚于 $V^n_2$。
\end{example*}

\begin{example*}{3}
闭曲面 $M^2$ 可以通过考察其上的特殊 Morse 函数轻易分类。
% [原文 p.177]
把这一方法在 $h$-配边定理的基础上推广，可得到某几类单连通流形的分类，例如只在中间维数 $[\frac n2]$ 有非零同调群的流形。Novikov（20 世纪 60 年代初）分类任意单连通流形的方法走的是另一条路线，尽管它确实用到了 $h$-配边定理（见下面 §4）。在任一连通闭 3 维流形 $M^3$ 上存在 Morse 函数 $f$，它有 1 个极大点、1 个极小点、一族指标为 1 的临界点（全在同一个水平面 $f=c_1$ 上）和一族指标为 2 的临界点（全在另一个水平面 $f=c_2$ 上）。（可以证明指标为 1 的临界点个数等于指标为 2 的个数；把这个数记作 $g$。）沿中间水平面 $f=c$（$c_1<c<c_2$）切开流形，得到两个子流形
\[
  M^3_-=\{x\in M\mid f(x)\le c\},\qquad
  M^3_+=\{x\in M\mid f(x)\ge c\},
\]
它们以曲面 $M^2_g$ 为公共边界：
\[
  M^2_g=\{x\in M\mid f(x)=c\}.
\]
不难看出，若 $M^3$ 可定向，则 $M_-$、$M_+$ 各自同胚于实心把手体：即把按标准方式嵌入 $\R^3$ 的、带 $g$ 个把手的球面 $M^2_g$ 的内部填满所得的 $\R^3$ 区域。由此可知，3 维流形 $M^3$ 的结构完全由这两个（相同的）实心把手体 $M_-$、$M_+$ 沿边界粘合的方式决定，即由一个微分同胚 $h:M^2_g\to M^2_g$ 决定。（3 维流形的这种分解称为 \emph{Heegaard 分解}。）虽然微分同胚 $M^2_g\to M^2_g$ 的同伦类与拓扑类已经分类（Nielsen，20 世纪 30 年代），但这并不给出（闭连通）3 维流形的分类，因为不同的微分同胚 $M^2_g\to M^2_g$ 可能给出同胚的流形 $M^3$。当 $g=1$ 时，具有亏格 1 的 Heegaard 分解的 3 维流形只有 $S^3$、$S^2\times S^1$ 以及某些透镜空间。
\end{example*}

\begin{example}{4}
设 $g(z)$ 是复维数为 $n$ 的复流形 $M^n$ 上处于一般位置的全纯函数，则 $g(z)$ 的实部 $f(z)=\mathrm{Re}\,g(z)$ 是 Morse 函数，其所有临界点的指标都是 $n$。特别地，这对由射影代数簇 $M^n\subset\CP^n$ 删去超平面截面 $\CP^{n-1}\cap M^n=W^{n-1}$ 所得的流形也成立：在剩余的仿射部分 $M^n\setminus W^{n-1}$ 上存在许多处于一般位置的处处全纯函数 $g$，它们的实部便是所有临界点指标都为 $n$ 的 Morse 函数。由此可以断定：每个映射 $K\to M^n$（其中 $K$ 是维数 $\le n-1$ 的 $CW$-复形）都可同伦地形变到超平面截面上。对同调群而言，这意味着包含同态
\[
  H_j(W^{n-1})\longrightarrow H_j(M^n)
\]
当 $j<n-1$ 时为同构，当 $j=n-1$ 时为满射。
% [Tutor Note: 此处 $M^n$ 的上标 $n$ 指复维数，作为实流形其维数是 $2n$；超平面截面 $W^{n-1}$ 的余维为 2（实维数 $2n-2$），不要把 $n-1$ 误当成余维数。]

% [原文 p.178]
由 Lefschetz 在 20 世纪 20 年代发展的射影代数簇的同调论，包括闭链的相交、代数闭链的理论以及超平面截面的特殊作用，将在综览复代数几何的相应卷册中详细考察。这里我们只限于考虑亚纯映射
\[
  f:M^n\longrightarrow\CP^1\cong S^2
\]
（即有极点的复解析函数）的若干特殊性质。这种函数的临界点伴随着一个漂亮的拓扑现象。设 $x_j$ 是临界点（$\mathrm{grad}\,f|_{x_j}=0$，$d^2f|_{x_j}$ 非退化），设 $f(x_j)=c$。在由方程 $f(x)=c+\eps$ 定义的每个非奇异水平集（或「纤维」）$F_\eps$ 上（$|\eps|$ 充分小），存在一个「消失闭链」$z_0\in H_{n-1}(F_\eps)$，所谓消失是指当 $\eps\to0$ 时该闭链趋于 0。让纤维 $F_\eps$ 沿闭路 $|\eps|=\mathrm{const}$ 绕行，就得到「单值性（monodromy）」映射
\[
  A:F_\eps\longrightarrow F_\eps.
\]
它在同调群上诱导的自映射 $A_j:H_j(F_\eps)\to H_j(F_\eps)$ 具有形式
\[
  A_j=1\quad\text{当 } j\ne n-1,
\]
\[
  A_{n-1}(z)=z\pm(z_0\cdot z)z_0,\qquad z,z_0\in H_{n-1}(F_\eps).
\]
（注意纤维 $F_\eps$ 是复维数为 $n-1$ 的复流形。）
% [Tutor Note: 所谓「消失」是指 $\eps\to0$ 时闭链 $z_0$ 在 $F_\eps$ 中收缩为 0 点；公式中 $\pm$ 的符号取决于绕行方向，并与临界点附近的局部模型有关。]

与代数函数的退化临界点相联系的单值性变换，长期以来（包括最近十年）是代数几何中深入研究的对象，参与者有 Deligne 等许多人。最有趣的情形之一是：用一族单参数超平面截割簇 $M^n\subset\CP^N$，得到这种「带奇点的纤维丛」，其纤维 $F$ 全都是超平面截面。从同调的角度看，在临界值 $c_j=g(x_j)$ 之上带奇异纤维的丛，借助球面 $S^2=\CP^1$ 上「Leray 层」$\mathcal F$ 中带系数的同调来描述，其中 $\mathcal F_U=H_*(g^{-1}(U))$。（注意这里对小圆盘 $U$，区域 $p^{-1}(U)$ 可缩到单一纤维上。）对只有非退化临界点的亚纯函数 $f:M^n\to\CP^1$ 而言，$f$ 的同调不变量只有各纤维 $F$ 的同调群，以及临界点的集合连同与之相联系的消失闭链（见上）。事实上这种情形早已被 20 世纪 20 年代的 Picard--Lefschetz 理论涵盖，只是严格的证明稍后才给出。
% [Tutor Note: 原书 p.178 此句作 $p^{-1}(U)$，而本段函数记作 $g$（$\mathcal F_U=H_*(g^{-1}(U))$）；上下文中的逆像应为 $g^{-1}(U)$，原书的 $p$ 疑为笔误，此处照原书排印。]
\end{example}

% [原文 p.179]
处理流形的实或复上同调理论还有一条完全不同的途径——用张量，即用反对称张量场（只带下指标），也就是微分形式，连同前文已引入的、作用在这种形式上的微分算子 $d$；这条途径由 Cartan 在 1930 年前后开创。某些重要类型的 1-形式与 2-形式连同它们的一些更深刻的拓扑性质与代数拓扑性质（包括应用），早先已由 Poincaré 结合自守形式理论和他对 Hamilton 系统的研究考察过；然而完整的系统理论归功于 Cartan。

考虑代数复形 $\Lambda^*(M^n)$，把 $M^n$ 上光滑的、整体定义的微分形式 $\Omega_k$ 作为它的元素，在每个坐标卡 $V_\alpha\subset M^n$ 上用局部坐标表为
\[
  \Omega_k=\sum_{j_1<\cdots<j_k}T_{j_1\cdots j_k}\,dx_\alpha^{j_1}\wedge\cdots\wedge dx_\alpha^{j_k}.
\]
这里 $k$ 是形式的次数，$k$-形式 $\Omega_k$ 合起来构成 $\Lambda^k(M^n)$。微分算子 $d:\Lambda^k(M^n)\to\Lambda^{k+1}(M^n)$ 局部地由
\begin{equation}
  d\Omega_k=\sum_{j_1<\cdots<j_k}\frac{\partial T_{j_1\cdots j_k}}{\partial x_\alpha^j}
  \,dx_\alpha^j\wedge dx_\alpha^{j_1}\wedge\cdots\wedge dx_\alpha^{j_k},\qquad
  dx_\alpha^j\wedge dx_\alpha^s=-dx_\alpha^s\wedge dx_\alpha^j \tag{2.14}
\end{equation}
定义。

实向量空间 $\Lambda^k(M^n)$（$k=0,\ldots,n$）连同微分 $d$ 确定代数复形 $\Lambda^*(M^n)$：
\begin{equation}
  0\longrightarrow\Lambda^0(M^n)\stackrel{d}{\longrightarrow}\Lambda^1(M^n)
  \longrightarrow\cdots\stackrel{d}{\longrightarrow}\Lambda^n(M^n)\longrightarrow0. \tag{2.15}
\end{equation}
注意 $\Lambda^k(M^n)$ 也可定义为 $k$-形式（$C^\infty$ 类）芽的层，而序列 (2.15) 给出层的正合列（见下文）。流形 $M^n$ 的上同调群现在由代数复形 (2.15) 的同调来定义：
\[
  H^k_\Lambda(M^n;\R)=\Ker d/\Image d,\qquad
  \Image d\subset\Ker d\subset\Lambda^k(M^n),
\]
它们的直和构成分次环
\[
  H^*_\Lambda(M^n)=\sum_{k\ge0}H^k_\Lambda(M^n;\R),
\]
其中乘法运算是形式的外积。「Poincaré 引理」断言：在 $M^n$ 每点的某邻域内，每个闭形式 $\Omega\in\Lambda^k(M^n)$（$k>0$，即 $d\Omega=0$）都是恰当的（即 $\Omega=d\Omega'$）。
% [Tutor Note: 这里的 $H^*_\Lambda$ 就是 de Rham 上同调；它与奇异上同调的相等正是 de Rham 定理。序列 (2.15) 局部正合而整体不正合，$\Ker d/\Image d$ 正是在度量整体障碍。]

对紧单连通 Lie 群以及某些齐性空间，即协变常曲率的单连通对称空间（也就是 Riemann 曲率的协变导数恒为零的空间），上同调环 $H^*_\Lambda(M^n)$ 容易算出。

% [原文 p.180]
在 Lie 群情形 $M^n=G$ 下，每个类 $z\in H^*_\Lambda(M^n)$ 都含有唯一的双不变 $k$-形式，反之每个这样的 $k$-形式都属于某个上同调类；而在对称空间情形 $M^n=G'/G'_0$ 下，每个类则含有唯一的 $G'$-不变 $k$-形式（反之亦然）。注意在后一情形，Lie 群 $G'$ 的 Lie 代数 $\mathfrak{g}'$ 总带有 $\Z/2$-分次：
\begin{align}
  \mathfrak{g}'&=\mathfrak{g}'_0+\mathfrak{g}'_1,\notag\\
  [\mathfrak{g}'_0,\mathfrak{g}'_0]&\subset\mathfrak{g}'_0,\quad
  [\mathfrak{g}'_0,\mathfrak{g}'_1]\subset\mathfrak{g}'_1,\quad
  [\mathfrak{g}'_1,\mathfrak{g}'_1]\subset\mathfrak{g}'_0. \tag{2.16}
\end{align}

赋有 Riemann 度量的 Lie 群 $G$ 可看作对称空间：取群 $G'=G\times G$ 作用于 $G$，作用由
\[
  (g_1,g_2):g\longmapsto g_1gg_2^{-1}
\]
给出。于是 $M^n=G\cong G'/G$，且
\begin{equation}
  \mathfrak{g}'=\mathfrak{g}'_0+\mathfrak{g}'_1,\qquad\text{其中}\quad
  \mathfrak{g}'_0=\{(x,x)\},\quad \mathfrak{g}'_1=\{(x,-x)\}. \tag{2.17}
\end{equation}

事实上对紧 Lie 群 $G$ 的任一齐性空间，上同调环的计算都归结为复形 $\Lambda^*_{inv}(M^n)$ 的计算，该复形由 $G$-不变形式构成；不过算子 $d$ 限制到 $\Lambda^*_{inv}$ 上，其作用可能相当于零算子。

20 世纪 30 年代 de Rham 证明：这些上同调群（即用微分形式定义的上同调群）自然地同构于通常以 $\R$ 为系数的单纯上同调群。现在我们用层论概述这一结果的现代证明。考虑流形 $M^n$ 上的下列层：
$\Lambda^k$ = $k$-形式（$C^\infty$ 类）芽的层；
$Z^k$ = 闭 $k$-形式芽的层。
（于是 $Z^0$ 是常数层。）有层的正合列
\begin{equation}
  0\longrightarrow Z^k\stackrel{i}{\longrightarrow}\Lambda^k\stackrel{d}{\longrightarrow}Z^{k+1}\longrightarrow0 \tag{2.18}
\end{equation}
（其中 $i$ 是包含映射），如第 3 章 §4 所述，它给出正合上同调列
\begin{equation}
  \cdots\longrightarrow H^j(Z^k)\stackrel{i_*}{\longrightarrow}H^j(\Lambda^k)
  \stackrel{d_*}{\longrightarrow}H^j(Z^{k+1})\stackrel{\partial^*}{\longrightarrow}
  H^{j+1}(Z^k)\longrightarrow\cdots, \tag{2.19}
\end{equation}
（这里 $H^j(\mathcal F)$ 表示以层 $\mathcal F$ 为系数的第 $j$ 个上同调群，即群 $H^j(M^n;\mathcal F)$。）利用 $\Lambda^k$（$h\ge0$）的截面只要求为 $C^\infty$ 类而不必解析这一事实，不难证明
% [Tutor Note: 原书 p.180 此处把变量印作 $h$（$\Lambda^k$（$h\ge0$）），而全节记号是 $k$（$\Lambda^k$、$Z^k$）；$h$ 疑为笔误，此处照原书排印。]
\begin{equation}
  H^j(\Lambda^k)=0\qquad\text{当 } j>0. \tag{2.20}
\end{equation}

% [原文 p.181]
还注意 $H^0(\Lambda^k)=\Lambda^k(M^n)$，因为按定义 $H^0(\Lambda^k)$ 就是整个 $M^n$ 上定义的 $C^\infty$ 类 $k$-形式的群。由 (2.20) 及序列 (2.19) 的正合性可知，当 $j>0$ 时同态
\[
  \partial^*:H^j(Z^{k+1})\longrightarrow H^{j+1}(Z^k)
\]
都是同构。于是我们有一串同构
\begin{equation}
  H^1(Z^k)\stackrel{\partial^*}{\longrightarrow}H^2(Z^{k-1})\stackrel{\partial^*}{\longrightarrow}
  \cdots\stackrel{\partial^*}{\longrightarrow}H^{k+1}(Z^0), \tag{2.21}
\end{equation}
其中
\begin{equation}
  H^{k+1}(Z^0)\cong H^{k+1}(M^n;\R), \tag{2.22}
\end{equation}
即通常的实上同调群，因为 $Z^0$ 是常数层。由 (2.19) 与 (2.20) 还有正合列
\[
  H^0(\Lambda^k)\stackrel{d_*}{\longrightarrow}H^0(Z^{k+1})
  \stackrel{\partial^*}{\longrightarrow}H^1(Z^k)\longrightarrow0,
\]
其中 $H^0(\Lambda^k)$ 是 $M^n$ 上整体定义的 $C^\infty$ 类 $k$-形式的群（于是 $d_*=d$），$H^0(Z^{k+1})$ 是 $M^n$ 上闭 $(k+1)$-形式的群。由这个正合列连同 (2.20)、(2.21) 立刻得到
\[
  H^{k+1}(M^n;\R)\cong H^1(Z^k)\cong\Ker d/\Image d\cong H^{k+1}_\Lambda(M^n),
  \qquad \Ker d\subset\Lambda^k(M^n),
\]
这就是 de Rham 定理。（注意此论证中隐含使用了 Poincaré 引理来确定局部正合性，即得到层序列 (2.18) 的正合性。）

下面我们来看上同调群 $H^j_\Lambda(M^n)$ 的另一类不变性，即它们在流形光滑同伦
\[
  F:M^n_1\times I\longrightarrow M^n_2,\qquad I=[a,b]
\]
下的不变性。首先注意，对任一形式 $\omega\in\Lambda^k(M^n_2)$，拉回具有形式
\[
  F^*\omega=\phi^{(0)}_{i_1\cdots i_{k-1}}dx^{i_1}\wedge\cdots\wedge dx^{i_{k-1}}\wedge dt
  +\phi^{(1)}_{j_1\cdots j_k}dx^{j_1}\wedge\cdots\wedge dx^{j_k}.
\]
若在形式 $\omega\in\Lambda^k(M^n_2)$ 上定义算子 $D$：
\[
  D\omega=\int_a^bF^*\omega,\qquad D\omega\in\Lambda_{k-1}(M^n_1),
\]
其中积分是对 $t$ 作的，则 $F^*\omega$ 的上式中第二项对 $D\omega$ 的贡献为零，因为该项不含 $dt$。利用这一点容易验证
\[
  Dd\pm dD=f_a^*-f_b^*,
\]

% [原文 p.182]
其中 $f_a=F(\,\cdot\,,a)$，$f_b=F(\,\cdot\,,b)$。于是（正如第 3 章 §2 中由该处式 (2.17) 得出类似结论那样）对任一上闭链 $z$（$dz=0$），差 $f_a^*(z)-f_b^*(z)=dD(z)$，从而是上边缘。因此诱导同态
\[
  f_a^*,f_b^*:H^j_\Lambda(M^n_2)\longrightarrow H^j_\Lambda(M^n_1)
\]
（对所有 $j$）相合。

就同伦论而言，$CW$-复形与胞腔映射的范畴可换成流形与光滑映射的范畴：任一 $CW$-复形 $K$ 嵌入某 $\R^N$ 后，可用它在 $\R^N$ 中可缩到 $K$ 的邻域代替，而胞腔映射可任意接近地用光滑映射逼近，或用可经任意小同伦形变到它的光滑映射代替。由这一考虑以及刚才概要证明的 $H^j_\Lambda$ 的同伦不变性（连同它们满足第 3 章 §3 中所列的其他上同调公理这一事实），显然可推出 $H^j_\Lambda$ 与通常上同调群相合。事实上沿着这条路线建立「用流形上的形式定义的上同调论」与通常上同调论的等价并无困难；不过详尽严格的阐述似乎不值得，因为它不会给出比已知其他证明更短的证明。

可定向流形 $M^n$ 上存在由 $M^n$ 的 Riemann 度量 $g_{ij}$ 确定的自然算子，称为 \emph{对偶算子}
$*:\Lambda^k(M^n)\to\Lambda^{n-k}(M^n)$。它局部地由
\[
  *\omega=(*T)_{i_1\cdots i_{n-k}}dx^{i_1}\wedge\cdots\wedge dx^{i_{n-k}},
\]
\[
  (*T)_{i_1\cdots i_{n-k}}=\sqrt{\det(g_{ij})}\,
  g_{i_1j_1}\cdots g_{i_{n-k}j_{n-k}}\,
  \varepsilon^{j_1\cdots j_{n-k}j_{n-k+1}\cdots j_n}T_{j_{n-k+1}\cdots j_n},
\]
\[
  \omega=T_{j_1\cdots j_k}dx^{j_1}\wedge\cdots\wedge dx^{j_k}
\]
给出。（算子 $*$ 对可定向流形 $M^n$ 定义良好，因为对这种流形，局部表达式
\[
  d\sigma=\sqrt{\det(g_{ij})}\,dx_\alpha^1\wedge\cdots\wedge dx_\alpha^n
\]
确定一个定义良好的 $n$-形式。）对偶算子 $*$ 具有下列性质：
1) 在 $k$-形式上 $(*)^2=(-1)^{k(n-k)}$；
2) 算子 $\delta=*^{-1}d*:\Lambda^{k+1}(M^n)\to\Lambda^k(M^n)$ 关于由
\begin{equation}
  \langle\omega,\omega'\rangle=\int_{M^n}\omega\wedge*\omega' \tag{2.23}
\end{equation}
给出的纯量积是 $d$ 的形式伴随算子；即
\[
  \pm\langle d\omega,\omega'\rangle=\langle\omega,\delta\omega'\rangle.
\]
% [原文 p.183]
3) Laplace 算子 $\Delta=d\delta+\delta d=(d+\delta)^2$ 与 $d$、$\delta$ 都可交换（且是自伴的）。

对任一闭流形 $M^n$，把 $\Lambda^k(M^n)$ 看作关于上述纯量积 (2.23) 的 Hilbert 空间，有如下正交「Hodge 分解」：
\begin{equation}
  \Lambda^k(M^n)=\Image d\oplus\Image\delta\oplus\Ker\Delta, \tag{2.24}
\end{equation}
其中 $\Ker\Delta$ 由 $M^n$ 上的 \emph{调和形式} 组成。还有
\begin{equation}
  \begin{aligned}
    \Ker d&=\Image d\oplus\Ker\Delta,\\
    \Ker\delta&=\Image\delta\oplus\Ker\Delta,\qquad
    \Ker\Delta=\Ker d\cap\Ker\delta,
  \end{aligned} \tag{2.25}
\end{equation}
由此可知 $\Lambda^k(M^n)$ 中调和形式空间的维数等于第 $k$ 个 Betti 数 $b_k(M^n)$。可以证明，Poincaré 对偶同态由对偶算子 $*$ 在调和形式上的作用决定。
% [Tutor Note: (2.24)、(2.25) 中的直和都是正交直和；由 $\Ker d=\Image d\oplus\Ker\Delta$ 立即得 $\Ker d/\Image d\cong\Ker\Delta$，即每个上同调类有唯一的调和代表，这正是 Hodge 定理的核心。]

Laplace 算子 $\Delta$ 也可对单纯复形定义，事实上可对更一般的 $CW$-复形定义：分别用 $\partial^*$ 与 $\partial$ 代替 $d$ 与 $\delta$，于是此处的 Laplace 算子定义为 $\Delta=\partial\partial^*+\partial^*\partial$。作此理解后，Hodge 分解的公式 (2.24)、(2.25) 对任一 $CW$-复形 $K$ 都成立，其纯量积由取各胞腔构成正交基确定。

若 $CW$-复形 $K$ 不单连通，则如第 3 章 §5 所述，任一（酉）表示 $\rho:\pi_1(K)\to U_n$ 确定一个具有 $\Z[\pi_1]$-模结构的链复形 $C_*(K;\rho)$。若该链复形所有正维数的同调群都为零，则 Reidemeister 挠有如下公式（Singer，20 世纪 60 年代末）：
\begin{equation}
  2\log|R(K,\rho)|=\sum_q q(-1)^q\log\det\Delta|_{C_q(K;\rho)}. \tag{2.26}
\end{equation}

把这一公式推广到流形上微分形式的复形，并利用适当定义的、两个不同表示的相对行列式，Ray 与 Singer（20 世纪 70 年代初）得以对 Reidemeister 挠作解析处理，特别证明了该概念的解析版本与度量 $g_{ij}$ 无关。Ray--Singer 挠确实与通常组合地定义的 Reidemeister 挠相合这一点的严格验证，稍后才以一般形式给出（由 Müller 与 Cheeger 在 20 世纪 70 年代末给出）。

此处理应提一下 A.\,S.\,Schwarz 对 Ray--Singer 挠的一个有趣的量子论解释：他对一个简单的、用微分形式写出的规范不变表达式作泛函积分，而挠出现在按标准方式构造一个假想地破坏规范不变性的理论的过程中。\footnote{一份综述写于 1983--1985 年，1986 年以俄文发表。上面提到的 A.\,S.\,Schwarz 的工作完成于 1980 年前后，当时已为某些数学物理学家（量子场论专家）所熟知。从 20 世纪 80 年代末起，这些思想由 Witten 等人发展。这一理论现在称为「拓扑量子场论」，已成为现代拓扑中极为活跃的领域。}
% [原文 p.184]

微分形式在复维数为 $n$、带 Kähler 度量 $g_{k\bar j}$ 的 Kähler 流形 $M^n$ 的同调论中特别重要。回想（上一节末尾）该度量按定义是 Hermite 度量，在以 $z_\alpha^1,\ldots,z_\alpha^n$ 为坐标的卡 $V_\alpha$ 上局部地由
\[
  ds^2=g_{k\bar j}\,dz_\alpha^k\,d\bar z_\alpha^{\bar j}
\]
给出，并具有如下性质：与之相伴的实 2-形式
\[
  \Omega=\frac{\ii}{2}g_{k\bar j}\,dz_\alpha^k\wedge d\bar z_\alpha^{\bar j}
\]
是闭的：$d\Omega=0$。$k$-形式的空间 $\Lambda^k(M^n)$（以及 $M^n$ 上相伴的层 $\Lambda^k$）自然地分解为 $(p,q)$ 型形式的直和：
\begin{equation}
  \Lambda^k(M^n)=\sum_{p+q=n}\Lambda^{p,q}, \tag{2.27}
\end{equation}
其中 $(p,q)$ 型形式（局部地）定义为
\begin{equation}
  \sum f_{i_1\cdots i_p\bar j_1\cdots\bar j_q}\,
  dz_\alpha^{i_1}\wedge\cdots\wedge dz_\alpha^{i_p}\wedge
  d\bar z_\alpha^{\bar j_1}\wedge\cdots\wedge d\bar z_\alpha^{\bar j_q}. \tag{2.28}
\end{equation}
% [Tutor Note: 原书 p.184 式 (2.27) 的求和指标印作 p+q=n，但 $\Lambda^k$ 的分解按 $k$ 分次，正确应为 p+q=k；下文 (2.32) 处也用了 p+q=k，可对照。]

Kähler 流形 $M^n$ 上形式的外微分，在 0-形式（函数）上由
\begin{equation}
  d=d_c+\bar d_c,\qquad d_cf=\frac{\partial f}{\partial z}dz,\qquad
  \bar d_cf=\frac{\partial f}{\partial\bar z}d\bar z \tag{2.29}
\end{equation}
定义，然后基本按实形式的方式扩充到 $M^n$ 上次数大于 0 的形式。Kähler 度量用来定义形式对偶的算子
\begin{equation}
  \delta=*^{-1}d*,\qquad \delta_c=*^{-1}d_c*,\qquad
  \bar\delta_c=*^{-1}\bar d_c*, \tag{2.30}
\end{equation}
以及相应的 Laplace 算子
\begin{equation}
  \begin{aligned}
    \Delta_c&=\delta_cd_c+d_c\delta_c,\qquad \Delta_c=\Delta_c,\\
    \Delta&=d\delta+\delta d.
  \end{aligned} \tag{2.31}
\end{equation}
% [Tutor Note: 原书 p.184 式 (2.31) 第二式印作 Delta_c = Delta_c，左端与右端同一记号，疑漏印横线；按 Kähler 恒等式应为 bar Delta_c = bar delta_c bar d_c + bar d_c bar delta_c，且 Delta_c 与 bar Delta_c 相等。]

由「Hodge 引理」，只要度量是 Kähler 度量，实 Laplace 算子与复 Laplace 算子就只相差一个常数因子：

% [原文 p.185]
\[
  \Delta=\mathrm{const.}\times\Delta_c.
\]
由此可知，对闭流形 $M^n$，调和形式的空间、从而上同调群分解为直和：
\begin{equation}
  H^k(M;\C)=\sum_{p+q=k}H^{p,q}(M^n;\C), \tag{2.32}
\end{equation}
这与分解 (2.27) 相一致。这一分解当然依赖于度量。注意由该度量确定的形式 $\Omega$（见上）是 $(1,1)$ 型的：
\[
  [\Omega]\in H^{1,1}(M^n;\C),
\]
并且 $[\Omega^n]\ne0$，因为积分
\[
  \int_{M^n}\Omega^n
\]
是 $M^n$ 体积的非零倍数。

用 $M^n$ 上任一形式 $\omega$ 与形式 $\Omega$ 作外积的运算通常记作 $L$：
\[
  L(\omega)=\Omega\wedge\omega,
\]
对偶算子记作 $\Lambda$。\emph{本原形式} $\omega$ 定义为满足 $\Lambda\omega=0$ 的形式。每个形式都可由本原形式经若干由 $L$ 与 $\Lambda$ 给出的算子作用而得到。注意 Dolbeault 的结果：
\[
  H^{p,q}(M^n;\C)\cong H^p(M^n;\Lambda^q),
\]
这里 $\Lambda^q$ 表示 $M^n$ 上全纯 $q$-形式芽的层。若形式 $\Omega$ 表示一个整上同调类，则由 Kodaira 定理，流形 $M^n$ 可代数地嵌入 $\CP^N$，使得 $\Omega$ 的上同调类对偶于超平面截面（前面已提到它在代数簇同调论中的特殊作用）。还注意 Moishezon 的结果（20 世纪 60 年代中期）：若复流形上的亚纯函数空间的超越次数（即代数无关函数个数的最大值）等于其维数，则该流形（带未必为 Hodge 的 Kähler 度量）可代数地嵌入 $\CP^N$。

一个有用的观察是：代数闭链，即射影代数簇 $M^n$ 的子簇（可以带奇点）$N^{n-k}$，确定对偶于 $H^{k,k}$ 中元素的上同调类：
\begin{equation}
  D[N^{n-k}]\in H^{k,k}(M^n;\C)\cap H^{2k}(M^n;\Z). \tag{2.33}
\end{equation}
按 Hodge 猜想的一个加强形式，式 (2.33) 右端交集中的每个元素都可用代数闭链的有理线性组合实现。
% [原文 p.186]
这一猜想 $k=1$ 的情形早先已由 Lefschetz 以更强的形式解决：$H^{1,1}\cap H^2(M^n;\Z)$ 中每个类都可用代数子簇 $N^{n-1}\subset M^n$（所谓「除子」）的基本闭链实现（如式 (2.33) 那样）。对 $k>1$，Thom 的理论（见下文）表明：即使要把 $H_{2(n-k)}(M^n;\Z)$ 中的元素实现为 $M^n$ 中非奇异复簇（甚至拟复簇）$N^{n-k}$ 的像（在映射 $N^{n-k}\to M^n$ 之下），也会出现障碍。这里还注意 Hironaka 20 世纪 60 年代中期的结果：每个带奇点的簇都可经收缩由非奇异簇得到（从而是它的连续像）。

本综述不讨论 Grothendieck 在 20 世纪 50 年代末开创的那一套原始想法，即为特征 $p>0$ 的（代数闭）域上的射影代数簇建立同调论。这些想法依赖于对 $\CP^N$ 开区域（相对 Zariski 拓扑，其闭集恰为代数子簇）的有限层无分歧覆叠范畴的巧妙运用，涉及 Grothendieck 的「étale 拓扑」。Grothendieck 理论中与光滑自映射 $M\to M$ 的不动点个数的 Lefschetz 公式相对应的结论，应用于有限域代数闭包的 Galois 群，并沿着 Weil--Tate 纲领用同调代数，就给出模 $p$ 丢番图方程解个数的数论渐近公式，以及 Riemann $\zeta$ 函数的 mod $p$ 类比的有理性（后者比通常的 $\zeta$ 函数简单得多）。这些问题连同有关历史与后来的发展，无疑将在本系列中专论代数几何的卷册中处理。对我们这里重要的是：由整系数多项式定义的实簇的 Grothendieck 同调论，与由同一组多项式定义的复簇的通常同调论给出的信息相同；情形只在 mod $p$ 同调中更复杂，这里 $p>0$ 是所关心的域的特征。

现在我们回到通常的实光滑流形理论。同调论与同伦论最重要的应用领域之一，是研究 $M^n$ 上道路泛函的极值曲线的整体性质，边界条件分下列两种类型之一：
\begin{equation}
  S\{\gamma\}=\int_{t_0}^{t_1}L(x,\dot x)\,dt,\qquad
  \gamma=x(t),\quad x(t_0)=x_0,\quad x(t_1)=x_1, \tag{2.34}
\end{equation}
即在道路 $\gamma$ 的端点固定于 $x_0$、$x_1$ 的极值问题；
\begin{equation}
  S(\gamma)=\oint_\gamma L(x,\dot x)\,dt,\qquad \gamma(t+T)=\gamma(t), \tag{2.35}
\end{equation}
即周期变分问题。

% CHUNK ch04d DONE
